% SPDX-License-Identifier: Apache-2.0
\section{A Card's Blocks Through Memory}
\label{sec:x-sddma}

A card sends the blocks of a multi-block read back to back, and the
card in Section~\ref{sec:x-sd} was given a gap between them to let a
program empty the buffer through a register. A real card gives
nothing like that. \code{Sd} therefore has a way to memory: the two
engines of Section~\ref{sec:x-dma}, \code{LineStore} for a read and
\code{LineFetch} for a write, and two more registers, \code{dma}, where
in memory the blocks start, and \code{blocks}, how many (issue~912).

A command with a count of blocks runs them on its own. The 128-word
buffer becomes a ring between the card's lines and an engine. For a
read, each word goes to the store engine as soon as the lines have
filled it, so the buffer never holds more than a few words. For a
write, the fetch engine's words go into the ring while it has room,
and a block goes out only once all of it is in, since the card cannot
be kept waiting in the middle of one. The command is done when its
last block is over and its engine has finished, so a program waits on
\code{status} as it does for any command, and then sends \code{CMD12}
to stop the card.

The two engines share one way onto the bus, an \code{Arbiter2}, and
never move words at the same time. A command is a read or a write and
never both, and no command starts until the last one, engine and all,
has finished. The host says as much in a \code{check!} that the two
busy lines are never high together, and the run counts the cycles in
which they are, which must be none.

The run brings the card up on four lines, writes two blocks from
memory to blocks 4 and 5 with \code{CMD25}, and reads them back into
another part of memory with \code{CMD18}, the card sending its blocks
a clock apart. It checks that the card holds what memory held and that
memory holds it again. The host is lowered, and its netlist is
simulated against the run under nvc and Verilator.

\begin{figure}[htbp]
\centering
\input{sddma_timing.tex}
\caption{The read's first block ending and the second starting: a
word to the store engine every eight card clocks, \code{fill} never
above one, \code{left} down to one at the block's end, and the store
engine busy throughout.}
\label{fig:x-sddma}
\end{figure}

\lstinputlisting[caption={\code{ex\_sddma.rs}. Run with \code{bazel run //lib/examples:ex\_sddma}.},label={lst:x-sddma}]{lib/examples/ex_sddma.rs}

\lstinputlisting[style=ascii,caption={What \code{ex\_sddma} printed when the build ran it: the card up, the two blocks written and read back, and no cycle with both engines running.},label={lst:x-sddma-out}]{out_sddma.txt}
